\chapter{Distinction Pack: Master Examiner Artefacts}
\label{app:distinction_pack}

\noindent\textbf{Purpose.} This appendix consolidates the eighteen highest-leverage examiner artefacts in a single navigable pack. Each section answers one specific examiner question and is intentionally a single table or short matrix rather than narrative. The pack functions as the dissertation's viva survival kit: an examiner can read it in fifteen minutes and reconstruct the entire study's logic, evidence, contribution, and boundary conditions.

\vspace{0.5em}
\noindent\textbf{How to read this appendix.} Sections are ordered by examiner priority. The single most important section (V.1, Dissertation-in-One-Page) compresses the entire study into one table; V.2 (Master Research Alignment Matrix) demonstrates that every claim is logically connected; V.3--V.10 establish the contribution at theory, novelty, knowledge, framework, and societal levels; V.11--V.14 document quality, risk, and reproducibility safeguards; V.15--V.18 position the work for publication and future research.

\section{Dissertation-in-One-Page}
\label{app:distinction_one_page}

\noindent\textbf{Examiner question.} What is the entire study in one image?

\noindent The table below is the canonical one-page summary of the dissertation. If a reviewer reads only one artefact, this should be it.

{\tablestripes\tablefontsize
\needspace{22\baselineskip}
\begin{xltabular}{\textwidth}{@{} >{\raggedright\arraybackslash}p{4.5cm} L @{}}
\caption[Dissertation in one page]{Dissertation in one page.}\label{tab:dist_one_page} \\
\toprule
\thead{Element} & \thead{One-sentence summary} \\
\midrule
\endfirsthead
\endhead
\bottomrule
\endlastfoot
Problem & Healthcare AI for ASD screening lacks an operational governance model that integrates explainability, trust, and adoption readiness. \\
\addlinespace
Gap & Governance is discussed in principles, not measured as an organisational capability that drives adoption. \\
\addlinespace
Objective & Validate whether governance maturity influences explainability, trust, and adoption readiness in B2C-primary, clinician-supported ASD screening. \\
\addlinespace
Theory & Extended Technology Acceptance Model (TAM) with governance maturity as an antecedent operating through explainability and trust. \\
\addlinespace
Method & Mixed methods: PLS-SEM on primary caregiver survey data plus secondary EEG benchmark validation, with cross-source triangulation. \\
\addlinespace
Sample & B2C caregivers (primary survey) plus public EEG benchmark datasets (secondary validation). \\
\addlinespace
Framework & RGAIG --- Responsible Generative AI Governance Framework with governance maturity, explainability, trust, and adoption readiness constructs. \\
\addlinespace
Headline finding & Governance Maturity $\rightarrow$ Explainability $\rightarrow$ Trust $\rightarrow$ Adoption Readiness; mediated path is stronger than the direct path. \\
\addlinespace
Contribution & Governance-centered adoption model for responsible healthcare AI, operationalised as a measurable framework with control catalogue. \\
\addlinespace
Impact & A deployable governance maturity assessment for clinical AI programmes; a publishable extension of TAM; a policy-readable responsible AI playbook. \\
\end{xltabular}}

\vspace{0.3em}\noindent\textit{Note.} Each row of the one-page summary is supported by the chapter and appendix references listed in subsequent sections of this Distinction Pack.

\section{Master Research Alignment Matrix}
\label{app:distinction_alignment}

\noindent\textbf{Examiner question.} How do I know everything aligns?

\noindent The matrix below traces every problem to its gap, objective, research question, hypothesis, method, finding, and contribution. End-to-end alignment is the strongest single signal of a doctoral-level study.

{\tablestripes\tablefontsize
\needspace{24\baselineskip}
\begin{xltabular}{\textwidth}{@{} p{1.0cm} p{1.7cm} p{2.0cm} p{1.4cm} p{1.0cm} p{1.7cm} L p{2.4cm} @{}}
\caption[Master research alignment matrix]{Master research alignment matrix.}\label{tab:dist_alignment} \\
\toprule
\thead{Prob.} & \thead{Gap} & \thead{Objective} & \thead{RQ} & \thead{H} & \thead{Method} & \thead{Finding} & \thead{Contribution} \\
\midrule
\endfirsthead
\endhead
\bottomrule
\endlastfoot
P1 & G1: governance not measured & O1: operationalise governance maturity & RQ1 & H1 & PLS-SEM on B2C survey & Governance maturity is a valid measurable construct & Governance Maturity construct (theory + scale) \\
\addlinespace
P2 & G2: explainability not linked to governance & O2: test governance$\rightarrow$explainability & RQ2 & H2 & PLS-SEM path & Governance significantly improves perceived explainability & Empirical link from governance to explainability \\
\addlinespace
P3 & G3: trust formation under-specified & O3: test explainability$\rightarrow$trust & RQ3 & H3 & PLS-SEM mediation & Explainability significantly improves trust & Confirmed mediator chain \\
\addlinespace
P4 & G4: adoption antecedents incomplete & O4: test trust$\rightarrow$adoption & RQ4 & H4 & PLS-SEM path & Trust significantly improves adoption readiness & Trust validated as proximal antecedent \\
\addlinespace
P5 & G5: no integrated governance$\rightarrow$adoption model & O5: integrate RGAIG framework & RQ5 & H5 & Cross-source triangulation & Mediated path is stronger than direct path & RGAIG framework validation \\
\end{xltabular}}

\vspace{0.3em}\noindent\textit{Note.} Each row's evidence is in Chapter~4; each contribution is elaborated in Chapter~5.

\section{Novelty Matrix}
\label{app:distinction_novelty}

\noindent\textbf{Examiner question.} What is genuinely new?

{\tablestripes\tablefontsize
\needspace{18\baselineskip}
\begin{xltabular}{\textwidth}{@{} >{\raggedright\arraybackslash}p{5.5cm} L @{}}
\caption[Novelty matrix]{Novelty matrix: existing research versus the present contribution.}\label{tab:dist_novelty} \\
\toprule
\thead{Existing research} & \thead{Present contribution} \\
\midrule
\endfirsthead
\endhead
\bottomrule
\endlastfoot
Governance discussed as principle & Governance measured as organisational capability \\
\addlinespace
Trust discussed as outcome & Trust validated as a mediating mechanism \\
\addlinespace
Explainable AI studied in isolation & Explainability integrated into a governance$\rightarrow$adoption pathway \\
\addlinespace
Adoption models without governance & Governance-extended adoption model (extended TAM) \\
\addlinespace
Responsible AI as principles statement & Operationalised RGAIG framework with measurable constructs \\
\addlinespace
XAI as technology-centric output & XAI as a governance-integration mechanism for stakeholder trust \\
\end{xltabular}}

\vspace{0.3em}\noindent\textit{Note.} Each row is one of the dimensions on which the dissertation extends prior research; the full lineage is in V.4 Theory Contribution Matrix.


\section{Theory Contribution Matrix}
\label{app:distinction_theory_contribution}

\noindent\textbf{Examiner question.} How does this advance theoretical knowledge?

{\tablestripes\tablefontsize
\needspace{18\baselineskip}
\begin{xltabular}{\textwidth}{@{} >{\raggedright\arraybackslash}p{4.5cm} L @{}}
\caption[Theory contribution matrix]{Theory contribution matrix --- existing theory and the present extension.}\label{tab:dist_theory_contribution} \\
\toprule
\thead{Existing theory} & \thead{Present contribution} \\
\midrule
\endfirsthead
\endhead
\bottomrule
\endlastfoot
Technology Acceptance Model (TAM) & Extended with governance maturity as an antecedent operating through explainability and trust \\
\addlinespace
Trust Theory & Repositioned with governance as a structural antecedent rather than an outcome \\
\addlinespace
Responsible AI principles & Operationalised into measurable constructs with empirical validation \\
\addlinespace
Socio-Technical Systems & Empirical evidence that joint optimisation requires explicit governance measurement \\
\end{xltabular}}

\vspace{0.3em}\noindent\textit{Note.} The four propositions (P1--P4) above are the journal-ready theoretical contribution.


\vspace{0.5em}\noindent\textbf{Formal propositions.} The theoretical contribution can be expressed as four propositions suitable for publication:
\begin{itemize}\setlength{\itemsep}{0pt}
\item[\textbf{P1.}] Governance maturity positively influences explainability.
\item[\textbf{P2.}] Explainability positively influences stakeholder trust.
\item[\textbf{P3.}] Trust positively influences adoption readiness.
\item[\textbf{P4.}] Governance maturity indirectly influences adoption through the mediated chain explainability$\rightarrow$trust.
\end{itemize}

\section{Knowledge Contribution Matrix}
\label{app:distinction_knowledge_contribution}

\noindent\textbf{Examiner question.} What new knowledge is being added?

{\tablestripes\tablefontsize
\needspace{18\baselineskip}
\begin{xltabular}{\textwidth}{@{} >{\raggedright\arraybackslash}p{4.5cm} L @{}}
\caption[Knowledge contribution matrix]{Knowledge contribution matrix.}\label{tab:dist_knowledge_contribution} \\
\toprule
\thead{Knowledge gap} & \thead{New knowledge contributed} \\
\midrule
\endfirsthead
\endhead
\bottomrule
\endlastfoot
Governance not operationalised as a measured construct & Governance Maturity validated as a measurable antecedent of explainability \\
\addlinespace
Explainability disconnected from organisational governance & Explainability empirically linked to governance maturity and to trust \\
\addlinespace
Trust formation under-theorised in AI contexts & Trust validated as a mediating mechanism between explainability and adoption \\
\addlinespace
Adoption antecedents incomplete for healthcare AI & Adoption readiness validated as a proximal outcome of the governance-mediated chain \\
\addlinespace
Responsible AI principles without empirical operationalisation & RGAIG framework provides an empirically grounded operational model \\
\end{xltabular}}

\vspace{0.3em}\noindent\textit{Note.} Each row closes a specific gap identified in the Chapter~2 literature review.


\section{Contribution Traceability Matrix}
\label{app:distinction_contribution_traceability}

\noindent\textbf{Examiner question.} Where in the dissertation does each contribution land?

{\tablestripes\tablefontsize
\needspace{12\baselineskip}
\begin{xltabular}{\textwidth}{@{} >{\raggedright\arraybackslash}p{4.5cm} p{2.0cm} L @{}}
\caption[Contribution traceability matrix]{Contribution traceability matrix.}\label{tab:dist_contribution_traceability} \\
\toprule
\thead{Contribution} & \thead{Chapter} & \thead{Evidence} \\
\midrule
\endfirsthead
\endhead
\bottomrule
\endlastfoot
Extended TAM with governance antecedent & Ch.~2 + Ch.~4 & Theory development (Ch.~2) and PLS-SEM verdicts (Ch.~4) \\
\addlinespace
Governance Maturity construct & Ch.~2 + Ch.~4 & Construct definition (Ch.~2), measurement model (Ch.~4) \\
\addlinespace
RGAIG framework & Ch.~5 & Framework synthesis and validation discussion \\
\addlinespace
Adoption model (G$\rightarrow$E$\rightarrow$T$\rightarrow$A) & Ch.~4 & Structural model + hypothesis verdicts \\
\addlinespace
Practitioner roadmap & Ch.~5 & Recommendation matrix and executive decision framework \\
\addlinespace
Methodological contribution (mixed-methods governance measurement) & Ch.~3 & Triangulation design + measurement model \\
\end{xltabular}}

\vspace{0.3em}\noindent\textit{Note.} Every contribution is reachable from one or more chapter sections; bidirectional cross-references are in V.19.


\section{Claim-Evidence Matrix}
\label{app:distinction_claim_evidence}

\noindent\textbf{Examiner question.} What claim rests on what evidence?

{\tablestripes\tablefontsize
\needspace{12\baselineskip}
\begin{xltabular}{\textwidth}{@{} >{\raggedright\arraybackslash}p{5.0cm} L @{}}
\caption[Claim--evidence matrix]{Claim--evidence matrix.}\label{tab:dist_claim_evidence} \\
\toprule
\thead{Claim} & \thead{Evidence} \\
\midrule
\endfirsthead
\endhead
\bottomrule
\endlastfoot
Governance maturity matters for adoption & PLS-SEM path Governance$\rightarrow$Adoption (mediated and direct), Chapter~4 \\
\addlinespace
Explainability mediates governance into trust & PLS-SEM mediation test, Chapter~4 \\
\addlinespace
Trust mediates explainability into adoption & PLS-SEM mediation test, Chapter~4 \\
\addlinespace
RGAIG framework is operationally valid & Construct reliability/validity tests + framework evaluation matrix \\
\addlinespace
Mediated path beats direct path & Comparison of path coefficients, Chapter~4 \\
\addlinespace
Findings generalise beyond a single dataset & Secondary EEG benchmark triangulation, Chapter~3 + Chapter~4 \\
\end{xltabular}}

\vspace{0.3em}\noindent\textit{Note.} Each claim is supported by a chapter, a table, or a figure; the full evidence trail is in Chapter~4 Integrated Findings.


\section{Examiner Attack Matrix}
\label{app:distinction_attack_matrix}

\noindent\textbf{Examiner question.} How do I answer the toughest viva challenges?

{\tablestripes\tablefontsize
\needspace{20\baselineskip}
\begin{xltabular}{\textwidth}{@{} >{\raggedright\arraybackslash}p{4.8cm} >{\raggedright\arraybackslash}p{5.5cm} L @{}}
\caption[Examiner attack matrix]{Examiner attack matrix --- the toughest viva challenges and their responses.}\label{tab:dist_attack_matrix} \\
\toprule
\thead{Examiner challenge} & \thead{Response} & \thead{Evidence} \\
\midrule
\endfirsthead
\endhead
\bottomrule
\endlastfoot
Why a DBA and not a PhD? & The contribution is governance-centered adoption, a managerial question; technical novelty is out of scope. & Ch.~1 (Why DBA signature diagram) \\
\addlinespace
Why this sample size? & Sample meets PLS-SEM power requirements and matches comparable governance-and-adoption studies. & Ch.~3 (Sampling Strategy) \\
\addlinespace
Why TAM as the base theory? & TAM is the most validated adoption theory; the contribution is its extension with governance maturity. & Ch.~2 (Theory Integration) \\
\addlinespace
Why governance as the antecedent? & Literature gap explicitly identifies absence of operational governance in adoption models. & Ch.~2 (Research Gap) \\
\addlinespace
Why ASD as the context? & ASD screening is the test bed; the governance findings generalise to other healthcare AI contexts. & Ch.~1 (Scope IN/OUT) \\
\addlinespace
Why EEG when AI accuracy is already established? & EEG is technical feasibility evidence, not the primary contribution. & Ch.~3 (Why EEG Exists statement) \\
\addlinespace
Why mixed methods? & Quantitative establishes statistical validity; qualitative validates contextual mechanism. & Ch.~3 (Mixed-Methods Justification) \\
\addlinespace
What if results don't replicate? & Triangulation across primary and secondary evidence reduces single-source risk; cross-source convergence supports robustness. & Ch.~3 (Triangulation Strategy) \\
\addlinespace
Is governance just a buzzword? & Governance is operationalised as a measurable construct with reliability and validity statistics. & Ch.~4 (Measurement Model) \\
\addlinespace
Why is RGAIG novel? & It integrates four existing traditions into a single operational framework with measurable constructs. & V.3 Novelty Matrix \\
\end{xltabular}}

\vspace{0.3em}\noindent\textit{Note.} Each challenge--response--evidence row functions as a viva FAQ; consolidated with Appendix~S2 Examiner Defense Pack.


\section{Research Quality Matrix}
\label{app:distinction_quality_matrix}

\noindent\textbf{Examiner question.} What quality safeguards are in place?

{\tablestripes\tablefontsize
\needspace{18\baselineskip}
\begin{xltabular}{\textwidth}{@{} >{\raggedright\arraybackslash}p{4.5cm} L @{}}
\caption[Research quality matrix]{Research quality matrix.}\label{tab:dist_quality_matrix} \\
\toprule
\thead{Quality criterion} & \thead{Evidence} \\
\midrule
\endfirsthead
\endhead
\bottomrule
\endlastfoot
Reliability & Cronbach's alpha and composite reliability per construct (Chapter~4 measurement model) \\
\addlinespace
Convergent validity & Average Variance Extracted (AVE) per construct (Chapter~4) \\
\addlinespace
Discriminant validity & HTMT ratios and Fornell--Larcker criterion (Chapter~4) \\
\addlinespace
Common method bias (CMB) & Harman's single-factor test and full collinearity VIF assessment (Chapter~3 + Chapter~4) \\
\addlinespace
Non-response bias & Early-versus-late respondent comparison (Chapter~3) \\
\addlinespace
Triangulation & Cross-source convergence between primary survey and secondary EEG benchmark (Chapter~3) \\
\addlinespace
Ethics & Informed consent, anonymisation, human oversight, ethics committee approval (Chapter~3 + Appendix~F) \\
\end{xltabular}}

\vspace{0.3em}\noindent\textit{Note.} Each criterion is reported with its statistical or procedural test in Chapter~3 + Chapter~4.


\section{Risk Register}
\label{app:distinction_risk_register}

\noindent\textbf{Examiner question.} What can go wrong and how is it mitigated?

{\tablestripes\tablefontsize
\needspace{12\baselineskip}
\begin{xltabular}{\textwidth}{@{} >{\raggedright\arraybackslash}p{4.2cm} >{\raggedright\arraybackslash}p{1.7cm} L @{}}
\caption[Risk register]{Risk register --- residual risks and mitigations.}\label{tab:dist_risk_register} \\
\toprule
\thead{Risk} & \thead{Impact} & \thead{Mitigation} \\
\midrule
\endfirsthead
\endhead
\bottomrule
\endlastfoot
Sample size limits subgroup analysis & Medium & Bootstrap resampling + power analysis documented \\
\addlinespace
Cross-sectional design limits causal inference & High & Mediation tests + longitudinal study proposed in future research \\
\addlinespace
Self-report bias in survey instrument & Medium & Validated scales + procedural remedies for CMB \\
\addlinespace
Generalisability limited to healthcare context & Medium & Explicit scope boundary + multi-context validation in future research \\
\addlinespace
Governance perceptions may shift over time & Medium & Cross-source triangulation reduces snapshot risk \\
\addlinespace
EEG technical evidence is benchmark-bound & Low & Framed as feasibility evidence, not primary contribution \\
\end{xltabular}}

\vspace{0.3em}\noindent\textit{Note.} Residual risks are restated alongside their mitigations in Chapter~5 limitations and Appendix~U Negative Evidence.


\section{Stakeholder Impact Matrix}
\label{app:distinction_stakeholder_matrix}

\noindent\textbf{Examiner question.} Who benefits and how?

{\tablestripes\tablefontsize
\needspace{12\baselineskip}
\begin{xltabular}{\textwidth}{@{} >{\raggedright\arraybackslash}p{3.8cm} L @{}}
\caption[Stakeholder impact matrix]{Stakeholder impact matrix.}\label{tab:dist_stakeholder_matrix} \\
\toprule
\thead{Stakeholder} & \thead{Benefit} \\
\midrule
\endfirsthead
\endhead
\bottomrule
\endlastfoot
Healthcare executives & Deployable governance maturity assessment for AI deployment decisions \\
\addlinespace
Governance boards & Operationalised control catalogue and audit-ready evidence trail \\
\addlinespace
Clinicians & Explainability-aware decision support that respects clinician oversight \\
\addlinespace
AI/ML teams & Governance integration patterns for explainability and trust by design \\
\addlinespace
Regulators & Empirically grounded responsible AI framework supporting policy formation \\
\addlinespace
Caregivers (B2C primary users) & Earlier screening support with transparent rationale and clinician backup \\
\addlinespace
Patients (children with suspected ASD) & Earlier identification and faster pathway to intervention \\
\end{xltabular}}

\vspace{0.3em}\noindent\textit{Note.} Each stakeholder benefit maps to a recommendation in V.13 Recommendation Matrix.


\section{Framework Evaluation Matrix}
\label{app:distinction_framework_evaluation}

\noindent\textbf{Examiner question.} How good is the RGAIG framework against established framework criteria?

{\tablestripes\tablefontsize
\needspace{12\baselineskip}
\begin{xltabular}{\textwidth}{@{} >{\raggedright\arraybackslash}p{4.0cm} L @{}}
\caption[Framework evaluation matrix]{Framework evaluation matrix (RGAIG).}\label{tab:dist_framework_evaluation} \\
\toprule
\thead{Criterion} & \thead{RGAIG result} \\
\midrule
\endfirsthead
\endhead
\bottomrule
\endlastfoot
Completeness & Covers governance, explainability, trust, adoption, and maturity layers --- complete \\
\addlinespace
Consistency & Constructs are mutually exclusive and collectively exhaustive within the adoption pathway \\
\addlinespace
Practicality & Each construct has measurable indicators usable in audit and assessment \\
\addlinespace
Scalability & Maturity layer enables organisational application from ad-hoc to optimised levels \\
\addlinespace
Governance coverage & Includes structural, procedural, and outcome controls \\
\addlinespace
Adoption coverage & Connects governance to perception (explainability, trust) and intent (adoption readiness) \\
\addlinespace
Empirical validation & Construct reliability and validity supported by Chapter~4 measurement model \\
\end{xltabular}}

\vspace{0.3em}\noindent\textit{Note.} Empirical validation evidence is reported in Chapter~4 measurement model and structural model sections.


\section{Recommendation Matrix}
\label{app:distinction_recommendation_matrix}

\noindent\textbf{Examiner question.} What should each audience do?

{\tablestripes\tablefontsize
\needspace{12\baselineskip}
\begin{xltabular}{\textwidth}{@{} >{\raggedright\arraybackslash}p{4.0cm} L @{}}
\caption[Recommendation matrix]{Recommendation matrix --- audience-specific actions.}\label{tab:dist_recommendation_matrix} \\
\toprule
\thead{Audience} & \thead{Recommendation} \\
\midrule
\endfirsthead
\endhead
\bottomrule
\endlastfoot
Healthcare organisations & Conduct a governance maturity assessment before any clinical AI deployment \\
\addlinespace
C-suite executives & Treat governance as a strategic adoption lever, not as a post-deployment audit task \\
\addlinespace
AI/ML teams & Build explainability into the model lifecycle as a governance requirement, not a feature add-on \\
\addlinespace
Regulators & Adopt measurable governance criteria (RGAIG-style) in healthcare AI policy guidance \\
\addlinespace
Researchers & Treat governance maturity as a measurable antecedent in future adoption studies \\
\addlinespace
Clinician educators & Integrate explainability literacy into clinical AI training curricula \\
\end{xltabular}}

\vspace{0.3em}\noindent\textit{Note.} Each recommendation traces to a finding in Chapter~4 and a managerial implication in Chapter~5.


\section{Managerial, Policy, and Societal Implications}
\label{app:distinction_implications}

\noindent\textbf{Examiner question.} What are the implications across managerial, policy, and societal levels?

{\tablestripes\tablefontsize
\needspace{18\baselineskip}
\begin{xltabular}{\textwidth}{@{} >{\raggedright\arraybackslash}p{3.5cm} >{\raggedright\arraybackslash}p{4.0cm} L @{}}
\caption[Managerial, policy, and societal implications]{Managerial, policy, and societal implications.}\label{tab:dist_implications} \\
\toprule
\thead{Level} & \thead{Area} & \thead{Implication} \\
\midrule
\endfirsthead
\endhead
\bottomrule
\endlastfoot
Managerial & Governance & Governance maturity is the strategic lever; explainability and trust are downstream consequences \\
\addlinespace
Managerial & Adoption & Adoption resistance is often a governance gap mis-diagnosed as a trust gap \\
\addlinespace
Policy & Responsible AI & Operational governance criteria should supplement principle-based statements \\
\addlinespace
Policy & Healthcare AI & Regulatory guidance should mandate governance maturity assessment for high-stakes AI \\
\addlinespace
Policy & Transparency & Explainability requirements should be set at the governance layer, not only the model layer \\
\addlinespace
Societal & Healthcare equity & Governance maturity reduces the risk of inequitable AI deployment across populations \\
\addlinespace
Societal & Trust in AI & Empirical evidence linking governance to trust strengthens the public-trust case for responsible AI \\
\addlinespace
Societal & Responsible AI ecosystem & The dissertation contributes a reusable governance model beyond healthcare AI \\
\end{xltabular}}

\vspace{0.3em}\noindent\textit{Note.} The three levels (managerial, policy, societal) are restated as dedicated sections in Chapter~5.


\section{Publication Roadmap}
\label{app:distinction_publication_roadmap}

\noindent\textbf{Examiner question.} What publications can come from this work?

{\tablestripes\tablefontsize
\needspace{12\baselineskip}
\begin{xltabular}{\textwidth}{@{} >{\raggedright\arraybackslash}p{5.0cm} L @{}}
\caption[Publication roadmap]{Publication roadmap --- candidate papers from the dissertation.}\label{tab:dist_publication_roadmap} \\
\toprule
\thead{Candidate paper} & \thead{Source chapters and journal area} \\
\midrule
\endfirsthead
\endhead
\bottomrule
\endlastfoot
Governance Maturity as an Antecedent of AI Adoption & Ch.~2 + Ch.~4; information systems / technology management journals \\
\addlinespace
RGAIG Framework for Responsible Healthcare AI & Ch.~5; healthcare informatics / responsible AI journals \\
\addlinespace
Trust as a Mediator in AI Adoption: Empirical Evidence & Ch.~4; trust research / behavioural IS journals \\
\addlinespace
Governance-Mediated Adoption Model: Extending TAM & Ch.~2 + Ch.~4; technology adoption / IS theory journals \\
\addlinespace
Operationalising Responsible AI for Healthcare: A Governance Lens & Ch.~5; AI policy / digital health policy journals \\
\addlinespace
Mixed-Methods Validation of Governance Constructs in AI Adoption & Ch.~3 + Ch.~4; research methodology / mixed methods journals \\
\end{xltabular}}

\vspace{0.3em}\noindent\textit{Note.} Each candidate paper draws on the chapters and appendices indicated in the source column.


\section{Research Legacy and Future Research Program}
\label{app:distinction_research_legacy}

\noindent\textbf{Examiner question.} What is the research program after this dissertation?

\noindent The figure below positions the dissertation as Phase~1 of a multi-phase responsible-AI research program. The chain shows how the governance-centered model generalises beyond healthcare to GenAI, agentic AI, and physical AI environments.

\begin{figure}[!htbp]
\centering
\begin{tikzpicture}[
  node distance=0.55cm,
  fbox/.style={rectangle, draw=ggunavy, fill=ggunavy!8, rounded corners=4pt, very thick,
    minimum width=8.2cm, minimum height=0.85cm, align=center, font=\fignodefont},
  fboxgold/.style={rectangle, draw=ggunavy, fill=ggugold!25, rounded corners=4pt, very thick,
    minimum width=8.2cm, minimum height=0.95cm, align=center, font=\fignodefont\bfseries},
  farr/.style={-{Stealth[length=3mm, width=2mm]}, very thick, ggunavy}
]
\node[font=\figtitlefont, ggunavy] (title) at (0,0.7) {Research legacy --- five-phase program};
\node[fbox, below=0.85cm of title] (n0) {Phase 1: Healthcare AI governance (this dissertation, ASD screening)};
\node[fbox, below=of n0] (n1) {Phase 2: Multi-site and cross-country validation};
\node[fbox, below=of n1] (n2) {Phase 3: Generative AI governance (GenAI-specific maturity)};
\node[fbox, below=of n2] (n3) {Phase 4: Agentic AI governance (multi-agent systems)};
\node[fbox, below=of n3] (n4) {Phase 5: Physical AI governance (robotics, IoT, embodied AI)};
\node[fboxgold, below=of n4] (n5) {Future trustworthy AI ecosystems};
\draw[farr] (n0) -- (n1);
\draw[farr] (n1) -- (n2);
\draw[farr] (n2) -- (n3);
\draw[farr] (n3) -- (n4);
\draw[farr] (n4) -- (n5);
\end{tikzpicture}
\caption[Research legacy --- five-phase program]{Research legacy --- five-phase program.}
\label{fig:dist_research_legacy}
\end{figure}

\vspace{0.3em}\noindent\textit{Note.} Phase~1 is the dissertation; Phases~2--5 form the post-doctoral research agenda enabled by the RGAIG framework.

\section{Reproducibility Pack}
\label{app:distinction_reproducibility}

\noindent\textbf{Examiner question.} Can another researcher reproduce this study?

{\tablestripes\tablefontsize
\needspace{12\baselineskip}
\begin{xltabular}{\textwidth}{@{} >{\raggedright\arraybackslash}p{5.0cm} L @{}}
\caption[Reproducibility pack]{Reproducibility pack --- assets available for replication.}\label{tab:dist_reproducibility} \\
\toprule
\thead{Asset} & \thead{Availability and location} \\
\midrule
\endfirsthead
\endhead
\bottomrule
\endlastfoot
Survey instrument & Full questionnaire in Appendix~F with construct--item mapping \\
\addlinespace
Dataset description & Sample profile, sampling frame, and response statistics in Chapter~3 \\
\addlinespace
Construct operationalisation & Construct--item mapping in Appendix~F, measurement model in Chapter~4 \\
\addlinespace
Analysis procedure & PLS-SEM steps documented in Chapter~3 methodology and Chapter~4 results \\
\addlinespace
SEM specification & Path model and indicator weights documented in Chapter~4 \\
\addlinespace
RGAIG framework & Complete framework in Appendix~D and Chapter~5 contribution discussion \\
\addlinespace
Secondary EEG benchmark & Dataset references and preprocessing steps in Appendix~I \\
\addlinespace
Drill methodology & Validation drill protocols in Appendix~Q \\
\addlinespace
Pipeline reference code & Code organisation reference in Appendix~R \\
\end{xltabular}}

\vspace{0.3em}\noindent\textit{Note.} Each asset is available either in this dissertation's appendices or in the supplementary materials referenced in Chapter~3.


\section{Final Master Scorecard}
\label{app:distinction_master_scorecard}

\noindent\textbf{Examiner question.} How does the dissertation score against the ten distinction-level dimensions?

{\tablestripes\tablefontsize
\needspace{18\baselineskip}
\begin{xltabular}{\textwidth}{@{} >{\raggedright\arraybackslash}p{4.0cm} >{\raggedright\arraybackslash}p{2.0cm} L @{}}
\caption[Final master scorecard]{Final master scorecard --- target distinction-level state.}\label{tab:dist_master_scorecard} \\
\toprule
\thead{Dimension} & \thead{Target} & \thead{Evidence in this dissertation} \\
\midrule
\endfirsthead
\endhead
\bottomrule
\endlastfoot
Originality & 10/10 & V.3 Novelty Matrix \\
\addlinespace
Theory & 10/10 & V.4 Theory Contribution Matrix + four propositions \\
\addlinespace
Methodology & 10/10 & Ch.~3 mixed-methods design + V.9 Research Quality Matrix \\
\addlinespace
Evidence & 10/10 & Ch.~4 PLS-SEM + V.7 Claim--Evidence Matrix \\
\addlinespace
Contribution & 10/10 & V.5 Knowledge Contribution + V.6 Contribution Traceability \\
\addlinespace
Governance / framework & 10/10 & V.12 Framework Evaluation Matrix (RGAIG) \\
\addlinespace
Practical impact & 10/10 & V.11 Stakeholder Impact Matrix + V.13 Recommendation Matrix \\
\addlinespace
Defense readiness & 10/10 & V.8 Examiner Attack Matrix + Appendix~S2 Defense Pack \\
\addlinespace
Publication potential & 10/10 & V.16 Publication Roadmap (six candidate papers) \\
\addlinespace
Legacy value & 10/10 & V.17 Research Legacy (five-phase research program) \\
\end{xltabular}}

\vspace{0.3em}\noindent\textit{Note.} Evidence column points to the specific Appendix~V section or chapter section that supports the score.


\section{Cross-Reference Map --- Appendix V $\leftrightarrow$ Chapter Sections}
\label{app:distinction_crossref_map}

\noindent\textbf{Examiner question.} For each Appendix~V section, where is the corresponding in-chapter section?

\noindent The table below provides the canonical bidirectional cross-reference between each Distinction Pack section and its in-chapter counterpart. An examiner can therefore read in either direction: from chapter narrative to appendix evidence pack, or from appendix pack to chapter context.

{\tablestripes\tablefontsize
\needspace{20\baselineskip}
\begin{xltabular}{\textwidth}{@{} >{\raggedright\arraybackslash}p{4.0cm} >{\raggedright\arraybackslash}p{4.5cm} L @{}}
\caption[Distinction Pack cross-reference map]{Distinction Pack cross-reference map --- Appendix V $\leftrightarrow$ chapter sections.}\label{tab:dist_crossref_map} \\
\toprule
\thead{Appendix V section} & \thead{In-chapter counterpart} & \thead{Cross-reference} \\
\midrule
\endfirsthead
\endhead
\bottomrule
\endlastfoot
V.1 Dissertation-in-One-Page & Front-matter Executive Summary & \autoref{tab:exec_summary_one_page} \\
\addlinespace
V.2 Master Research Alignment Matrix & Chapter alignment tables (Ch.~4 hypothesis summary + Ch.~5 contribution traceability) & \autoref{tab:ch4_hypothesis_summary} + \autoref{tab:ch4_rq_mapping} \\
\addlinespace
V.3 Novelty Matrix & Ch.~1 Novelty Statement & \autoref{sec:ch1_novelty_statement} \\
\addlinespace
V.4 Theory Contribution Matrix + propositions P1--P4 & Ch.~5 Theory Contribution dedicated section & \autoref{sec:ch5_theory_contribution_dedicated} \\
\addlinespace
V.5 Knowledge Contribution Matrix & Ch.~2 + Ch.~5 Knowledge Contribution dedicated sections & \autoref{sec:ch2_knowledge_contribution_dedicated} + \autoref{sec:ch5_knowledge_contribution_dedicated} \\
\addlinespace
V.6 Contribution Traceability Matrix & Ch.~1 Contribution Statement & \autoref{sec:ch1_contribution_statement} \\
\addlinespace
V.7 Claim-Evidence Matrix & Ch.~4 Integrated Findings Summary & \autoref{sec:ch4_integrated_findings} \\
\addlinespace
V.8 Examiner Attack Matrix & Ch.~1 Why DBA + per-chapter Examiner Questions landing pages & \autoref{sec:ch1_why_dba_dedicated} + \autoref{sec:ch1_three_questions} \\
\addlinespace
V.9 Research Quality Matrix & Ch.~3 CMB + Non-Response Bias controls & \autoref{sec:ch3_cmb_controls} + \autoref{sec:ch3_nonresponse_bias} \\
\addlinespace
V.10 Risk Register & Appendix~U Negative Evidence + Ch.~3 methodology limitations & \autoref{app:negative_evidence} \\
\addlinespace
V.11 Stakeholder Impact Matrix & Ch.~5 Societal Contribution & \autoref{sec:ch5_societal_contribution} \\
\addlinespace
V.12 Framework Evaluation Matrix (RGAIG) & Ch.~4 Integrated Findings + Ch.~5 framework discussion & \autoref{sec:ch4_integrated_findings} \\
\addlinespace
V.13 Recommendation Matrix & Ch.~5 Recommendation Matrix dedicated section & \autoref{sec:ch5_recommendation_matrix_dedicated} \\
\addlinespace
V.14 Managerial/Policy/Societal Implications & Ch.~5 Managerial + Policy + Societal dedicated sections & \autoref{sec:ch5_policy_implications} + \autoref{sec:ch5_societal_contribution} \\
\addlinespace
V.15 Publication Roadmap & Future research agenda (Ch.~5) & \autoref{sec:ch5_template_compliance} \\
\addlinespace
V.16 Research Legacy (5-phase program) & Future research roadmap & \autoref{fig:dist_research_legacy} \\
\addlinespace
V.17 Reproducibility Pack & Appendix~F Survey Instrument + Appendix~Q Drill Methodology & \autoref{app:negative_evidence} \\
\addlinespace
V.18 Final Master Scorecard & Per-chapter Examiner Summaries & \autoref{sec:ch1_examiner_summary} \\
\end{xltabular}}

\vspace{0.3em}\noindent\textit{Note.} The cross-reference map is the single navigation index for the Distinction Pack. Every appendix-V section is reachable from at least one chapter section, and every dedicated chapter section is reachable from at least one appendix-V section.

\vspace{0.5em}\noindent\textbf{Closing statement.} The dissertation argues that healthcare AI adoption is mediated by governance maturity through the chain explainability$\rightarrow$trust$\rightarrow$adoption. The RGAIG framework operationalises this chain into measurable constructs, an evaluation matrix, and a five-phase research program that generalises beyond ASD screening to the broader responsible AI agenda. Every claim in the dissertation can be traced through this Distinction Pack to its supporting chapter, table, and figure.
